% ================================================================
% SECTION VIII: CONCLUSION
% ================================================================

We measured six production DA protocols (Polkadot ELVES, Ethereum EIP-4844/PeerDAS, Celestia, Espresso Tiramisu, NEAR and Avail) with one methodology, using 90~days of mainnet data covering over 14~million blocks.

\subsection{Summary of Findings}

Every protocol's mean utilization was below 40\%.
The networks are demand-constrained, so until they approach saturation, cross-protocol throughput comparisons measure adoption.

Median DA cost spans more than eight orders of magnitude, from $\sim$\$3$\times$10$^{-9}$/MiB (Espresso, whose 1~wei/byte fee is a Mainnet~0 placeholder) to $\sim$\$0.53/MiB (NEAR).
On Ethereum, the blob fee left the 1-wei floor because of Fusaka's EIP-7918 reserve price, before the BPO forks added capacity.
Annualizing cost over each protocol's retention window shows that a protocol that is cheap at spot can be expensive when data must stay available.

No protocol slashes for pure data withholding.
Polkadot gates finality on availability bitfields, Celestia uses bad-encoding fraud proofs, Ethereum uses peer scoring, Espresso relies on committee reputation, NEAR has an unused challenge mechanism, and Avail slashes only for equivocation.
Attributing withholding, an omission fault, to a specific actor is still unsolved.

No protocol guarantees permanent availability.
Protocol-defined retention ranges from 25~hours (Polkadot) to $\sim$18~days (Ethereum), Avail and Espresso leave retention to node operators, and only Ethereum's window comfortably exceeds a 7-day optimistic-rollup challenge period.

\subsection{Future Work}

\emph{Longitudinal extension}: repeating these measurements as utilization rises would show supply-constrained behaviour that the current data cannot.
\emph{Stress testing}: controlled synthetic load would measure protocol limits directly, which observational data cannot do.
\emph{Withholding enforcement}: incentive-compatible enforcement for omission faults needs formal analysis.
\emph{Cross-layer analysis}: measuring end-to-end rollup performance (confirmation time, cost, finality) for each DA backend would show how the choice affects rollups.
\emph{PeerDAS at scale}: PeerDAS activated during our window; its effect on node bandwidth and on further BPO increases is the next thing to measure.
